
\section{Reproducible numerics and scripts}
\label{sec:experiments}

This section provides optional reproducible scripts illustrating intermediate mechanisms discussed in the paper.
No closed-layer theorem depends on computation; the scripts are included to make the constructions auditable and testable.

\subsection{Reproducibility protocol}
\label{sec:repro_protocol}

All experiments are implemented in Python~3.10+.
The Poisson solver uses \texttt{numpy} if available (FFT mode); if not, it falls back to a pure-Python iterative periodic solver (slower, but dependency-free).
All other scripts use only the Python standard library.
Run the following from this paper directory:
\begin{itemize}
\item \path{python3 scripts/exp_hilbert_locality.py}
\item \path{python3 scripts/exp_golden_discrepancy_bound.py}
\item \path{python3 scripts/exp_abel_pole_barrier.py}
\item \path{python3 scripts/exp_poisson_fft.py}
\item \path{python3 scripts/exp_wigner_smith_kappa.py}
\item \path{python3 scripts/exp_address_family_sensitivity.py}
\item \path{python3 scripts/exp_kappa_redshift_toy.py}
\item \path{python3 scripts/exp_1p1d_routing_worked_example.py}
\end{itemize}

The scripts optionally write small \LaTeX{} row files into \path{sections/generated/}.
If those files are present, the tables below will include them.

\subsubsection*{Determinism and randomness}
All scripts are deterministic as written.
The only use of pseudorandomness is in the ``shuffled baseline'' of \path{scripts/exp_address_family_sensitivity.py}, and the RNG seeds are fixed in code (2D: \texttt{seed = 123456 + order}; 3D: \texttt{seed = 654321 + order}).
Thus, regenerating tables does not require passing external seeds.

\subsubsection*{Generated artifacts (what each script writes)}
Each script writes into \path{sections/generated/} using fixed filenames; tables include them when present.
{\small
\begin{itemize}
\item \textbf{Hilbert locality (Experiment A).}\\
\path{exp_hilbert_locality.py} writes \path{hilbert_locality_rows.tex}.\\
\emph{Defaults:} orders $n=1,\dots,8$.
\item \textbf{Address-family sensitivity (Experiments A', A'', A''', A3D).}\\
\path{exp_address_family_sensitivity.py} writes\\
\path{address_kappa_proxy_rows.tex},\\
\path{address_neighbor_model_sensitivity_rows.tex},\\
\path{address_kappa_proxy_fit_rows.tex},\\
\path{address_kappa_proxy_3d_rows.tex},\\
and (where used) \path{address_neighbor_separation_rows.tex}.\\
\emph{Defaults:} 2D orders $n=1,\dots,8$ with Manhattan and Chebyshev neighborhoods.
3D orders $n=1,\dots,5$.
Shuffled baseline seeds fixed as stated above.
\item \textbf{Golden discrepancy (Experiment B).}
\path{exp_golden_discrepancy_bound.py} writes \path{golden_discrepancy_rows.tex}.
\emph{Defaults:} $\alpha=1/\varphi$, $x_0=0.1$, and $N$ in \{100, 300, 1000, 3000, 10000, 30000\}.
\item \textbf{Abel pole barrier (Experiment C).}
\path{exp_abel_pole_barrier.py} writes \path{abel_barrier_rows.tex}.
\emph{Defaults:} $\gamma=1$ and $T_{\max}=5000$.
$\beta$ in \{0.5, 0.6\}, with $r$ sampled on a fixed grid around the threshold.
\item \textbf{Poisson FFT (Experiments D and D').}
\path{exp_poisson_fft.py} writes \path{poisson_rows.tex} and \path{poisson_scaling_rows.tex}.
\emph{Defaults:} FFT mode uses $N=64$ and $r_{\max}=12$ (fallback: $N=24$).
The scaling table uses $N$ in \{32, 48, 64\} (FFT) plus a modest Dirichlet-Jacobi comparison at $N=32$, with window band starting at $r=3$.
\item \textbf{Wigner--Smith interface (Experiment E).}\\
\path{exp_wigner_smith_kappa.py} writes \path{wigner_smith_rows.tex}.\\
\emph{Defaults:} toy Breit--Wigner with $E_0=1$ and $\gamma=0.2$.
Energy band $[0,2]$ sampled at $n=17$ points.
$\tau_0=1$, smoothing window $=1$ (off).
\item \textbf{Scan-chain redshift toy (Experiment F).}\\
\path{exp_kappa_redshift_toy.py} writes \path{kappa_redshift_toy_rows.tex}.\\
\emph{Defaults:} order $n=8$ and horizon $t_{\max}=2\times 10^8$ ticks.
Maps: Hilbert and Z-order.
\item \textbf{1+1D worked example (Appendix~\ref{app:worked_1p1d}).}
\path{exp_1p1d_routing_worked_example.py} writes
\path{1p1d_error_rows.tex},
\path{1p1d_scaling_fit_rows.tex},
and \path{1p1d_curve_rows.tex}.
\emph{Defaults:} weak-field target with $M=0.05$ on $r$ in $[1,8]$, and $\kappa_0=10^9$.
Orders $n=6,\dots,11$ for scaling (representative curve samples at order $n=9$).
\end{itemize}
}

\subsection{Experiment A: Hilbert address locality (2D)}
\label{sec:exp_hilbert_locality}

We verify the one-step locality property \eqref{eq:hilbert_locality} for the standard $2$D discrete Hilbert map up to order $n=8$.

\begin{table}[t]
\centering
\begin{tabular}{@{}rcc@{}}
\toprule
order $n$ & $\min_t \norm{H_n(t+1)-H_n(t)}_1$ & $\max_t \norm{H_n(t+1)-H_n(t)}_1$ \\
\midrule
\IfFileExists{sections/generated/hilbert_locality_rows.tex}{\input{sections/generated/hilbert_locality_rows.tex}}{(run \path{scripts/exp_hilbert_locality.py} to generate)}\\
\bottomrule
\end{tabular}
\caption{Hilbert address one-step locality check for orders $n=1,\dots,8$.}
\label{tab:hilbert_locality}
\end{table}

\subsection[Experiment A': address-family sensitivity on a scan chain (2D)]{Experiment A$'$: address-family sensitivity on a scan chain (2D)}
\label{sec:exp_address_sensitivity}

We quantify how the choice of address family changes scan-order separations of screen neighbors.
For an address map $A_n$ and a screen-neighbor edge $\{x,y\}$, define the separation $\Delta_{A_n}(x,y):=\abs{A_n^{-1}(x)-A_n^{-1}(y)}$ (Definition~\ref{def:neighbor_index_separation}).
On a scan chain, $\Delta_{A_n}$ lower bounds the depth required to enact neighbor interactions (Proposition~\ref{prop:scan_chain_separation_lower_bound}).
We report statistics of the induced local proxy $\widetilde\kappa_{A_n}(x)=\max_{y\sim x}\Delta_{A_n}(x,y)$ across all sites in $\Sigma_n$ for Hilbert vs.\ Morton/Z-order, together with a deterministic shuffled baseline.
Unless otherwise stated, the neighborhood relation $y\sim x$ here refers to the 4-neighbor (Manhattan) screen adjacency.

\begin{table}[t]
\centering
\begin{tabular}{@{}r l rrrrr@{}}
\toprule
order $n$ & map & mean & $p_{50}$ & $p_{90}$ & $p_{99}$ & max \\
\midrule
\IfFileExists{sections/generated/address_kappa_proxy_rows.tex}{1 & Hilbert & 2.00 & 1 & 3 & 3 & 3 \\
1 & Z-order & 2.00 & 2 & 2 & 2 & 2 \\
1 & Shuffled & 2.50 & 2 & 3 & 3 & 3 \\
2 & Hilbert & 5.25 & 3 & 11 & 13 & 13 \\
2 & Z-order & 4.25 & 3 & 6 & 6 & 6 \\
2 & Shuffled & 9.56 & 11 & 12 & 15 & 15 \\
3 & Hilbert & 12.94 & 7 & 43 & 53 & 53 \\
3 & Z-order & 10.31 & 6 & 22 & 22 & 22 \\
3 & Shuffled & 34.33 & 33 & 51 & 62 & 62 \\
4 & Hilbert & 30.25 & 11 & 75 & 211 & 213 \\
4 & Z-order & 24.08 & 11 & 86 & 86 & 86 \\
4 & Shuffled & 147.07 & 146 & 217 & 243 & 243 \\
5 & Hilbert & 66.92 & 13 & 179 & 819 & 853 \\
5 & Z-order & 53.64 & 22 & 171 & 342 & 342 \\
5 & Shuffled & 565.18 & 565 & 822 & 962 & 1010 \\
6 & Hilbert & 142.42 & 19 & 301 & 2869 & 3413 \\
6 & Z-order & 115.10 & 22 & 342 & 1366 & 1366 \\
6 & Shuffled & 2319.87 & 2309 & 3403 & 3856 & 4054 \\
7 & Hilbert & 295.62 & 19 & 341 & 5201 & 13653 \\
7 & Z-order & 240.49 & 22 & 342 & 5462 & 5462 \\
7 & Shuffled & 9228.63 & 9138 & 13536 & 15521 & 16341 \\
8 & Hilbert & 604.24 & 19 & 681 & 13133 & 54613 \\
8 & Z-order & 493.86 & 22 & 683 & 10923 & 21846 \\
8 & Shuffled & 37072.80 & 36807 & 54171 & 62123 & 65327 \\
}{(run \path{scripts/exp_address_family_sensitivity.py} to generate)}\\
\bottomrule
\end{tabular}
\caption{Address-family sensitivity on the 2D screen: statistics of the local scan-chain overhead proxy $\widetilde\kappa_{A_n}$ for Hilbert, Morton/Z-order, and a shuffled baseline.}
\label{tab:address_kappa_proxy}
\end{table}

\subsection[Experiment A'': neighborhood-model robustness (fixed order)]{Experiment A$''$: neighborhood-model robustness (fixed order)}
\label{sec:exp_address_neighbor_model}

To assess robustness of the proxy to adjacency conventions, we recompute the local proxy at a fixed order,
\[
\widetilde\kappa_{A_n}(x)=\max_{y\sim x}\Delta_{A_n}(x,y),
\]
using both the Manhattan 4-neighbor adjacency and the Chebyshev 8-neighbor adjacency (including diagonals).

\begin{table}[t]
\centering
\begin{tabular}{@{}l rrr rrr@{}}
\toprule
map & $p_{99}^{(4)}$ & $p_{99}^{(8)}$ & ratio &
max$^{(4)}$ & max$^{(8)}$ & ratio \\
\midrule
\IfFileExists{sections/generated/address_neighbor_model_sensitivity_rows.tex}{Hilbert & 13133 & 13136 & 1.000 & 54613 & 54613 & 1.000 \\
Z-order & 10923 & 10929 & 1.001 & 21846 & 32769 & 1.500 \\
Shuffled & 62123 & 62955 & 1.013 & 65327 & 65444 & 1.002 \\
}{(run \path{scripts/exp_address_family_sensitivity.py} to generate)}\\
\bottomrule
\end{tabular}
\caption{Sensitivity of the overhead proxy to the neighborhood model at a fixed resolution (script default: order $n=8$). Here $(4)$ denotes Manhattan 4-neighbors and $(8)$ denotes Chebyshev 8-neighbors.}
\label{tab:address_neighbor_model_sensitivity}
\end{table}

\subsection[Experiment A''': finite-size trend fit for high quantiles]{Experiment A$'''$: finite-size trend fit for high quantiles}
\label{sec:exp_address_fit}

To summarize finite-size scaling, we fit $\log_2 p_{99}$ and $\log_2(\max)$ as affine functions of $n$ over the computed range $n=1,\dots,8$ (least squares), and report the fitted slopes and $R^2$.

\begin{table}[t]
\centering
\begin{tabular}{@{}l r r r r@{}}
\toprule
map & slope for $\log_2 p_{99}$ & $R^2$ & slope for $\log_2(\max)$ & $R^2$ \\
\midrule
\IfFileExists{sections/generated/address_kappa_proxy_fit_rows.tex}{Hilbert & 1.751 & 0.987 & 2.015 & 1.000 \\
Z-order & 1.856 & 0.997 & 1.939 & 0.999 \\
Shuffled & 2.027 & 1.000 & 2.041 & 1.000 \\
}{(run \path{scripts/exp_address_family_sensitivity.py} to generate)}\\
\bottomrule
\end{tabular}
\caption{Finite-size trend fit for high quantiles of the Manhattan proxy $\widetilde\kappa_{A_n}$.}
\label{tab:address_kappa_proxy_fit}
\end{table}

\subsection[Experiment A3D: address sensitivity in three dimensions (scan-chain proxy)]{Experiment A$_{3D}$: address sensitivity in three dimensions (scan-chain proxy)}
\label{sec:exp_address_sensitivity_3d}

As a minimal dimensional robustness check, we repeat the same scan-chain proxy construction on a 3D screen lattice $\Sigma_n=\{0,\dots,2^n-1\}^3$ with 6-neighbor (Manhattan) adjacency, comparing Morton/Z-order against a shuffled baseline.
We report the local proxy $\widetilde\kappa_{A_n}(x)=\max_{y\sim x}\abs{A_n^{-1}(x)-A_n^{-1}(y)}$ statistics across all sites.

\begin{table}[t]
\centering
\begin{tabular}{@{}r l rrrrr@{}}
\toprule
order $n$ & map & mean & $p_{50}$ & $p_{90}$ & $p_{99}$ & max \\
\midrule
\IfFileExists{sections/generated/address_kappa_proxy_3d_rows.tex}{1 & Z-order & 4.00 & 4 & 4 & 4 & 4 \\
1 & Shuffled & 5.00 & 5 & 7 & 7 & 7 \\
2 & Z-order & 18.88 & 14 & 28 & 28 & 28 \\
2 & Shuffled & 37.38 & 37 & 51 & 59 & 59 \\
3 & Z-order & 92.83 & 55 & 220 & 220 & 220 \\
3 & Shuffled & 309.24 & 305 & 439 & 489 & 503 \\
4 & Z-order & 426.14 & 220 & 1756 & 1756 & 1756 \\
4 & Shuffled & 2487.73 & 2473 & 3499 & 3933 & 4069 \\
5 & Z-order & 1852.72 & 220 & 7022 & 14044 & 14044 \\
5 & Shuffled & 20068.47 & 19897 & 28014 & 31372 & 32656 \\
}{(run \path{scripts/exp_address_family_sensitivity.py} to generate)}\\
\bottomrule
\end{tabular}
\caption{3D scan-chain proxy statistics for Morton/Z-order and a shuffled baseline (orders $n=1,\dots,5$).}
\label{tab:address_kappa_proxy_3d}
\end{table}

\subsection{Experiment B: golden-branch star discrepancy and an explicit bound}
\label{sec:exp_golden_discrepancy}

We compute $\DNstar(P_N)$ for the golden scan and compare it against the explicit bound \eqref{eq:golden_bound}.

\begin{table}[t]
\centering
\begin{tabular}{@{}rccc@{}}
\toprule
$N$ & $\DNstar(P_N)$ & bound $2(2+\log_\varphi N)/N$ & ratio \\
\midrule
\IfFileExists{sections/generated/golden_discrepancy_rows.tex}{100 & 0.0207764 & 0.231399 & 0.090 \\
300 & 0.00575141 & 0.092353 & 0.062 \\
1000 & 0.00134748 & 0.0327098 & 0.041 \\
3000 & 0.000692067 & 0.0124253 & 0.056 \\
10000 & 0.000244267 & 0.00422798 & 0.058 \\
30000 & 8.00778e-05 & 0.00156153 & 0.051 \\
}{(run \path{scripts/exp_golden_discrepancy_bound.py} to generate)}\\
\bottomrule
\end{tabular}
\caption{Exact star discrepancy for the golden-branch Kronecker sequence and the explicit logarithmic bound.}
\label{tab:golden_discrepancy}
\end{table}

\subsection{Experiment C: Abel pole barrier toy model}
\label{sec:exp_abel_barrier}

We illustrate Lemma~\ref{lem:abel_pole_barrier} using a toy mode
$u_t=\e^{(\beta-\frac12)t}\cos(\gamma t)$.
When $\beta>\tfrac12$, the Abel weight must satisfy $r<\e^{-(\beta-\frac12)}$ to suppress exponential gain; otherwise the partial sums exhibit threshold blow-up.

\begin{table}[t]
\centering
\begin{tabular}{@{}ccc@{}}
\toprule
$\beta$ & $r$ & $\bigl|S(r;T_{\max})\bigr|$ \\
\midrule
\IfFileExists{sections/generated/abel_barrier_rows.tex}{0.500 & 0.8000 & 7.321034e-01 \\
0.500 & 0.8800 & 6.369817e-01 \\
0.500 & 0.9000 & 6.134388e-01 \\
0.500 & 0.9050 & 6.075852e-01 \\
0.500 & 0.9200 & 5.901151e-01 \\
0.500 & 0.9500 & 5.556554e-01 \\
0.500 & 0.9800 & 5.219656e-01 \\
0.600 & 0.8000 & 6.320969e-01 \\
0.600 & 0.8800 & 5.302520e-01 \\
0.600 & 0.9000 & 5.058303e-01 \\
0.600 & 0.9050 & 1.910336e+00 \\
0.600 & 0.9200 & 1.193817e+36 \\
0.600 & 0.9500 & 5.656057e+105 \\
0.600 & 0.9800 & 1.823025e+173 \\
}{(run \path{scripts/exp_abel_pole_barrier.py} to generate)}\\
\bottomrule
\end{tabular}
\caption{Toy Abel-threshold behavior at fixed $T_{\max}$.}
\label{tab:abel_barrier}
\end{table}

\subsection{Experiment D: Poisson phase potential via FFT}
\label{sec:exp_poisson}

We solve $-\Delta\Phi=4\pi\rho$ on a periodic grid and check that a localized source produces an approximate $1/r$ potential profile before periodic images dominate.

\begin{table}[t]
\centering
\begin{tabular}{@{}rcc@{}}
\toprule
$r$ & $\langle\Phi\rangle$ & $r\langle\Phi\rangle$ \\
\midrule
\IfFileExists{sections/generated/poisson_rows.tex}{1 & +0.704947 & +0.704947 \\
2 & +0.335675 & +0.671351 \\
3 & +0.203919 & +0.611756 \\
4 & +0.131873 & +0.527493 \\
5 & +0.082613 & +0.413063 \\
6 & +0.051132 & +0.306791 \\
7 & +0.031001 & +0.217009 \\
8 & +0.016342 & +0.130732 \\
9 & +0.004706 & +0.042356 \\
10 & -0.003518 & -0.035177 \\
}{(run \path{scripts/exp_poisson_fft.py} to generate)}\\
\bottomrule
\end{tabular}
\caption{Representative radial averages from the FFT Poisson solver (periodic box).}
\label{tab:poisson_rows}
\end{table}

\subsection[Experiment D': finite-size and source-width robustness of the 1/r window]{Experiment D$'$: finite-size and source-width robustness of the $1/r$ window}
\label{sec:exp_poisson_scaling}

To quantify how cleanly a periodic finite grid reproduces an intermediate $1/r$ regime, we measure the flatness of $r\langle\Phi\rangle$ over a fixed radius band and report the relative RMS deviation as a simple error proxy, for multiple grid sizes and for a point source vs.\ a small cube-smeared source.
For comparison, we also include a Dirichlet (zero boundary) iterative solve at a modest size.

\begin{table}[t]
\centering
\small
\begin{tabular}{@{}r l l l r r r@{}}
\toprule
$N$ & boundary/solver & source & window & count & mean($r\langle\Phi\rangle$) & rel.\ RMS \\
\midrule
\IfFileExists{sections/generated/poisson_scaling_rows.tex}{20 & periodic-Jacobi & cube-0 & 3--5 & 3 & +4.3128e-01 & 2.228e-01 \\
20 & periodic-Jacobi & cube-1 & 3--5 & 3 & +4.3283e-01 & 2.194e-01 \\
24 & periodic-Jacobi & cube-0 & 3--6 & 4 & +4.6478e-01 & 2.481e-01 \\
24 & periodic-Jacobi & cube-1 & 3--6 & 4 & +4.6574e-01 & 2.457e-01 \\
20 & Dirichlet-Jacobi & cube-0 & 3--5 & 3 & +6.0764e-01 & 1.104e-01 \\
20 & Dirichlet-Jacobi & cube-1 & 3--5 & 3 & +6.0714e-01 & 1.094e-01 \\
}{(run \path{scripts/exp_poisson_fft.py} to generate)}\\
\bottomrule
\end{tabular}
\caption{Finite-size and source-width robustness proxy for the $1/r$ window.}
\label{tab:poisson_scaling}
\end{table}

\subsection{Experiment E: Wigner--Smith time delay and \texorpdfstring{$\kappa_{\mathrm{WS}}(E)$}{kappaWS(E)}}
\label{sec:exp_wigner_smith}

To connect computational lapse to measurable delays, we include a minimal implementation of the Wigner--Smith time-delay matrix
\[
Q(E)=-\iu S(E)^\dagger \frac{\dd S}{\dd E},
\]
the total delay $\tau_{\mathrm{WS}}(E)=\Tr Q(E)$, and the dimensionless overhead proxy $\kappa_{\mathrm{WS}}(E)=\tau_{\mathrm{WS}}(E)/\tau_0$ for a chosen reference tick duration $\tau_0$.
The script \path{scripts/exp_wigner_smith_kappa.py} provides an interface and includes a 1-channel Breit--Wigner toy model; users may replace it with a physical model or data for the scattering matrix $S(E)$.

\begin{table}[t]
\centering
\begin{tabular}{@{}rcc@{}}
\toprule
$E$ & $\tau_{\mathrm{WS}}(E)$ & $\kappa_{\mathrm{WS}}(E)$ \\
\midrule
\IfFileExists{sections/generated/wigner_smith_rows.tex}{0.125 & 0.263063 & 0.263063 \\
0.250 & 0.358906 & 0.358906 \\
0.375 & 0.518752 & 0.518752 \\
0.500 & 0.815577 & 0.815577 \\
0.625 & 1.464887 & 1.464887 \\
0.750 & 3.313108 & 3.313108 \\
0.875 & 9.522320 & 9.522320 \\
1.000 & 14.336886 & 14.336886 \\
1.125 & 9.522320 & 9.522320 \\
1.250 & 3.313108 & 3.313108 \\
1.375 & 1.464887 & 1.464887 \\
1.500 & 0.815577 & 0.815577 \\
1.625 & 0.518752 & 0.518752 \\
1.750 & 0.358906 & 0.358906 \\
1.875 & 0.263063 & 0.263063 \\
}{(run \path{scripts/exp_wigner_smith_kappa.py} to generate)}\\
\bottomrule
\end{tabular}
\caption{Representative values from the Wigner--Smith toy model, providing an auditable interface for $\kappa_{\mathrm{WS}}(E)$.}
\label{tab:wigner_smith_rows}
\end{table}

\subsection{Experiment F: a scan-chain redshift toy from a computed \texorpdfstring{$\kappa(x)$}{kappa(x)} landscape}
\label{sec:exp_kappa_redshift_toy}

We demonstrate the redshift ratio \eqref{eq:redshift_ratio} on a fully discrete scan-chain toy model.
Fix a finite-resolution address map $A_n$ (Hilbert or Morton/Z-order) and place screen sites on a scan chain by $\pi_n(x)=A_n^{-1}(x)$.
Define a site-local overhead proxy by
\[
\kappa(x):=\widetilde\kappa_{A_n}(x)=\max_{y\sim x}\abs{A_n^{-1}(x)-A_n^{-1}(y)},
\]
as in Definition~\ref{def:neighbor_index_separation}.
Interpreting one local ``clock cycle'' at $x$ as requiring $\kappa(x)$ global ticks, the lapse dictionary predicts
\[
\frac{\dd\tau_1}{\dd\tau_2}=\frac{\kappa(x_2)}{\kappa(x_1)}.
\]
The script below selects representative low/median/high overhead sites and verifies the ratio by counting completed cycles over a long horizon.

\begin{table}[t]
\centering
\begin{tabular}{@{}r l l r l r rr@{}}
\toprule
order $n$ & map & $x_1$ & $\kappa(x_1)$ & $x_2$ & $\kappa(x_2)$ & predicted & measured \\
\midrule
\IfFileExists{sections/generated/kappa_redshift_toy_rows.tex}{8 & Hilbert & (0,255) & 1 & (2,7) & 19 & 19.000000 & 19.000001 \\
8 & Hilbert & (0,255) & 1 & (127,0) & 54613 & 54613.000000 & 54614.964500 \\
8 & Hilbert & (2,7) & 19 & (127,0) & 54613 & 2874.368421 & 2874.471600 \\
8 & Z-order & (0,0) & 2 & (0,3) & 22 & 11.000000 & 11.000000 \\
8 & Z-order & (0,0) & 2 & (0,127) & 21846 & 10923.000000 & 10924.186148 \\
8 & Z-order & (0,3) & 22 & (0,127) & 21846 & 993.000000 & 993.107822 \\
}{(run \path{scripts/exp_kappa_redshift_toy.py} to generate)}\\
\bottomrule
\end{tabular}
\caption{Scan-chain toy redshift verification from a computed overhead proxy $\kappa(x)$ (Hilbert vs Morton/Z-order).}
\label{tab:kappa_redshift_toy}
\end{table}


